\section{フッ化水素の物性と酸解離}

\subsection{令和元年度 問9}
\label{q:r01-reg-applied-09}

\noindent\textbf{原典：} sources/R01/17\_applied\_science.pdf、PDF物理3頁。

\textbf{公式正答：} 5（single）。

\par\textbf{検討未確定：} 純液体HFの誘電率と単独HFの金・白金非反応性について条件付き一次資料の全文確認が残る。発煙も濃度条件を要する。初回5と公式5の一致だけで全肢校正完了としない。

\paragraph{問題の要旨} フッ化水素の物性・ガラスへの反応・危険性と、酸解離の強さを区別する。

\Needspace{4\baselineskip}\paragraph{確認した条件・記述（独自要約）}\mbox{}\par

\noindent\textbf{確認範囲：} 問題幹、否定、全5肢、数式、数値、単位を原典画像で確認した。\par

\noindent\textbf{成立条件：} 他のハロゲン化水素との比較。純液体の誘電性と水溶液の酸解離は別。単独のHFと他の酸化剤を併用する系を区別する。\par

\paragraph{問われる核心} フッ化水素の物性・ガラスへの反応・危険性と、酸解離の強さを区別する。

\paragraph{必要知識} 他のハロゲン化水素との比較。純液体の誘電性と水溶液の酸解離は別。単独のHFと他の酸化剤を併用する系を区別する。

\paragraph{考え方・正答を選ぶ理由（査読待ち初稿）} HFは強い水素結合を形成し、他のハロゲン化水素より高い沸点を持つ。ガラスのケイ酸成分と反応する。高濃度水溶液は発煙し、弱酸であっても深刻な有害性を持つ。5は有害性を強酸であることに結び付けるため誤りである。

\Needspace{4\baselineskip}\paragraph{選択肢の検討（査読待ち初稿）}\mbox{}\par

\noindent\textbf{1：適切} 他のハロゲン化水素より沸点が高い。分子間水素結合による会合が寄与する。\par

\noindent\textbf{2：解釈未解決} 液体HFは大きな誘電率を持つ。定性的な候補として初回は適切としたが、温度条件付きの定量一次資料の照合が残る。\par

\noindent\textbf{3：条件付きで適切} ガラスと反応するが金・白金とは反応しない。ガラス腐食はNIOSHで確認できる。金・白金の非反応性を他の酸化剤との混合条件まで一般化しない。\par

\noindent\textbf{4：条件付きで適切} 水溶液は空気中で発煙する。NIOSHは濃度40パーセント超の発煙を説明する。濃度を指定しない全水溶液への全称とは読まない。\par

\noindent\textbf{5：不適切} 強い毒性を持つ強酸である。毒性・腐食性は強いが、通常の希薄水溶液における解離は弱酸の分類である。\par

\paragraph{再利用する考え方} 反応速度と平衡位置、全圧と分圧、濃度と活量を分ける。式の適用条件を確認し、原典の不正確な記述を科学的に正しい文へ置き換えない。

\paragraph{共通理論} \hyperref[theory:chem_eq:acid_strength]{第\ref*{theory:chem_eq:acid_strength}節}

\par 参考：\cite{purdue_acid_strength,niosh_hf}。

\par 難易度 2（校準初稿）、優先度未評価。

\section{Alの両性・不動態・塩と製錬を比較する。}

\subsection{平成26年度 問12}
\label{q:h26-reg-applied-12}

\noindent\textbf{原典：} sources/H26/17\_applied\_science.pdf、PDF物理5頁。

\textbf{公式正答：} 1（single）。

\par\textbf{検討未確定：} 濃硝酸の温度・濃度・表面条件、ミョウバン結晶化条件は未指定。原典の代表工程を全条件の保証にしない。

\paragraph{問題の要旨} Alの両性・不動態・塩と製錬を比較する。

\Needspace{4\baselineskip}\paragraph{確認した条件・記述（独自要約）}\mbox{}\par

\noindent\textbf{判断対象：} Alの両性・不動態・塩と製錬を比較する。\par

\noindent\textbf{原典確認：} 22原典頁画像で五候補・否定・化学式を照合。\par

\paragraph{問われる核心} Alの両性・不動態・塩と製錬を比較する。

\paragraph{必要知識} Alは両性で酸にも強塩基にも反応し、水溶液中の強塩基ではアルミン酸/水酸化錯イオン形成を伴うため1。

\paragraph{考え方・正答を選ぶ理由（査読待ち初稿）} Alは両性で酸にも強塩基にも反応し、水溶液中の強塩基ではアルミン酸/水酸化錯イオン形成を伴うため1。不動態は酸への全条件無反応を意味しない。ボーキサイト精製→アルミナ→氷晶石を含む融液の電解と、Al主合金・複塩の分類を別に整理する。

\Needspace{4\baselineskip}\paragraph{選択肢の検討（査読待ち初稿）}\mbox{}\par

\noindent\textbf{1：選ぶ候補} 酸と反応するが塩基とは反応しない。Alは強塩基とも反応し1が誤り。\par

\noindent\textbf{2：選ばない候補} 濃硝酸で酸化膜・不動態。通常の濃硝酸で表面保護が起きるという条件付き説明。\par

\noindent\textbf{3：選ばない候補} ジュラルミンはAl主成分。Alを基材とする合金分類。\par

\noindent\textbf{4：選ばない候補} Al硫酸塩/K硫酸塩からミョウバン。KAl(SO4)2・12H2Oという複塩。\par

\noindent\textbf{5：選ばない候補} ボーキサイト/アルミナ/氷晶石電解。アルミナを溶融塩に溶かして還元する。\par

\paragraph{再利用する考え方} 原子番号と周期・族、物性と反応性、酸強度と危険性、生成と除去、定義の分母を別々に確かめる。

\paragraph{共通理論} \hyperref[theory:inorganic:aluminum]{第\ref*{theory:inorganic:aluminum}節}

\par 参考：\cite{inorganic_al_ref}。

\par 難易度 2（校準初稿）、優先度未評価。

\section{レアメタルの元素と機器用途を対応させる。}

\subsection{令和4年度 問13}
\label{q:r04-reg-applied-13}

\noindent\textbf{原典：} sources/R04/17\_applied\_science.pdf、PDF物理4頁。

\textbf{公式正答：} 1（single）。

\par\textbf{検討未確定：} 機器は代表方式・出題時用途を扱う。全車両モータNd不可欠の全称、海洋鉱床の経済性・数量は未確定。

\paragraph{問題の要旨} レアメタルの元素と機器用途を対応させる。

\Needspace{4\baselineskip}\paragraph{確認した条件・記述（独自要約）}\mbox{}\par

\noindent\textbf{判断対象：} レアメタルの元素と機器用途を対応させる。\par

\noindent\textbf{原典確認：} 22原典頁画像で五候補・否定・化学式を照合。\par

\paragraph{問われる核心} レアメタルの元素と機器用途を対応させる。

\paragraph{必要知識} 代表的青色LEDはGaN/InGaNの窒化物半導体で、Asとする1が不適切。

\paragraph{考え方・正答を選ぶ理由（査読待ち初稿）} 代表的青色LEDはGaN/InGaNの窒化物半導体で、Asとする1が不適切。In透明電極、Nd磁石、Pt触媒は元素の電子構造と機器機能の対応。海洋資源の存在確認を商業採掘済みと読み替えず、Ndを全HV/EV方式に不可欠とする全称にも留保を残す。

\Needspace{4\baselineskip}\paragraph{選択肢の検討（査読待ち初稿）}\mbox{}\par

\noindent\textbf{1：選ぶ候補} 青LEDにAs。代表的青LEDはGaN/InGaN系で、Asの対応が不適切。\par

\noindent\textbf{2：選ばない候補} 透明電極にIn。ITO等の透明導電膜のIn。\par

\noindent\textbf{3：選ばない候補} 日本EEZに熱水/Coクラスト/RE泥。存在確認と経済的採掘・埋蔵量評価は別。\par

\noindent\textbf{4：選ばない候補} HV/EVモータにNd。Nd永久磁石モータの代表用途。全モータ方式に不可欠ではない。\par

\noindent\textbf{5：選ばない候補} 排ガス浄化触媒にPt。白金族金属の代表的触媒用途。\par

\paragraph{再利用する考え方} 原子番号と周期・族、物性と反応性、酸強度と危険性、生成と除去、定義の分母を別々に確かめる。

\paragraph{共通理論} \hyperref[theory:inorganic:materials]{第\ref*{theory:inorganic:materials}節}

\par 参考：\cite{inorganic_blueled,inorganic_rareuses,inorganic_searesources}。

\par 難易度 2（校準初稿）、優先度未評価。

\section{HFの酸化膜除去とPTFEの原料・加工を比較する。}

\subsection{令和元年度再試験 問9}
\label{q:r01-ret-applied-09}

\noindent\textbf{原典：} sources/R01\_retest/17\_applied\_science.pdf、PDF物理4頁。

\textbf{公式正答：} 3（single）。

\par\textbf{検討未確定：} PTFEの原料製造・全グレード加工の出題当時本文は未照合。初回と再結晶で融解温度が異なり得る。

\paragraph{問題の要旨} HFの酸化膜除去とPTFEの原料・加工を比較する。

\Needspace{4\baselineskip}\paragraph{確認した条件・記述（独自要約）}\mbox{}\par

\noindent\textbf{判断対象：} HFの酸化膜除去とPTFEの原料・加工を比較する。\par

\noindent\textbf{原典確認：} 22原典頁画像で五候補・否定・化学式を照合。\par

\paragraph{問われる核心} HFの酸化膜除去とPTFEの原料・加工を比較する。

\paragraph{必要知識} HFの半導体工程での用途は酸化膜の除去であり、生成とする3が逆。

\paragraph{考え方・正答を選ぶ理由（査読待ち初稿）} HFの半導体工程での用途は酸化膜の除去であり、生成とする3が逆。SiO2とFの反応による可溶性/揮発性フッ化物を使う。PTFEはTFE重合体で高い溶融粘度を持ち、特殊加工を使うが327℃を全熱履歴の普遍融点へ固定しない。

\Needspace{4\baselineskip}\paragraph{選択肢の検討（査読待ち初稿）}\mbox{}\par

\noindent\textbf{1：選ばない候補} FはPauling EN最大。指定尺度でFが最大。\par

\noindent\textbf{2：選ばない候補} HFはホタル石由来。CaF2からHFを製造。\par

\noindent\textbf{3：選ぶ候補} HFはウェハSiO2膜生成。HFはSiO2をエッチして除去するので生成とは逆。\par

\noindent\textbf{4：選ばない候補} TFE重合でPTFE。CF2=\allowbreak{}CF2の付加重合。\par

\noindent\textbf{5：選ばない候補} PTFE融点327℃・特殊成形。再結晶材の代表融点で、新生材などの熱履歴差を保持。\par

\paragraph{再利用する考え方} 原子番号と周期・族、物性と反応性、酸強度と危険性、生成と除去、定義の分母を別々に確かめる。

\paragraph{共通理論} \hyperref[theory:inorganic:fluorine]{第\ref*{theory:inorganic:fluorine}節}

\par 参考：\cite{inorganic_metalloid_ref,inorganic_hf_etch,inorganic_ptfe}。

\par 難易度 2（校準初稿）、優先度未評価。

\section{気体の水溶性・酸塩基・酸化還元を比較する。}

\subsection{平成23年度 問13}
\label{q:h23-reg-applied-13}

\noindent\textbf{原典：} sources/H23/17\_applied\_science.pdf、PDF物理5頁。

\textbf{公式正答：} 5（single）。

\par\textbf{検討未確定：} 最軽量は通常分子気体の同温同圧比較。NH3工業合成史の初めてという全称・全気体性状の原典当時全文は未照合。

\paragraph{問題の要旨} 気体の水溶性・酸塩基・酸化還元を比較する。

\Needspace{4\baselineskip}\paragraph{確認した条件・記述（独自要約）}\mbox{}\par

\noindent\textbf{判断対象：} 気体の水溶性・酸塩基・酸化還元を比較する。\par

\noindent\textbf{原典確認：} 22原典頁画像で五候補・否定・化学式を照合。\par

\paragraph{問われる核心} 気体の水溶性・酸塩基・酸化還元を比較する。

\paragraph{必要知識} H2SのSは−2で、例えばハロゲンへ電子を渡してSなどへ酸化される還元剤として扱う。

\paragraph{考え方・正答を選ぶ理由（査読待ち初稿）} H2SのSは−2で、例えばハロゲンへ電子を渡してSなどへ酸化される還元剤として扱う。5が誤り。硫化物の沈殿はpHと溶解度積の別機構である。CO2+Ca(OH)2→CaCO3+H2O、さらにCO2と水で炭酸水素イオンを含む溶液となる。臭い・毒性と酸化還元性を同一視しない。

\Needspace{4\baselineskip}\paragraph{選択肢の検討（査読待ち初稿）}\mbox{}\par

\noindent\textbf{1：選ばない候補} H2は最軽量で難溶。通常の中性分子気体の同温同圧比較。\par

\noindent\textbf{2：選ばない候補} Cl2は黄緑・酸化・漂白。水中HOCl等の酸化作用。\par

\noindent\textbf{3：選ばない候補} NH3は易溶アルカリ/Haber–Bosch。NH3の塩基性とN2/H2の合成。歴史の初めての範囲は留保。\par

\noindent\textbf{4：選ばない候補} CO2で石灰水白濁/過剰で溶解。CaCO3生成後に炭酸水素イオンを伴って溶解する。\par

\noindent\textbf{5：選ぶ候補} H2Sは腐卵臭・酸化性・硫化物沈殿。通常の反応比較ではH2Sは還元剤、酸化性という語が不適切。\par

\paragraph{再利用する考え方} 原子番号と周期・族、物性と反応性、酸強度と危険性、生成と除去、定義の分母を別々に確かめる。

\paragraph{共通理論} \hyperref[theory:inorganic:gases]{第\ref*{theory:inorganic:gases}節}

\par 参考：\cite{inorganic_halogen_ref,inorganic_hydrogen_ref}。

\par 難易度 2（校準初稿）、優先度未評価。

\section{ハロゲン単体の状態・酸化力とHFの酸強度を比較する。}

\subsection{平成28年度 問11}
\label{q:h28-reg-applied-11}

\noindent\textbf{原典：} sources/H28/17\_applied\_science.pdf、PDF物理4頁。

\textbf{公式正答：} 5（single）。

\par\textbf{検討未確定：} HFの酸強度は希薄水溶液の通常分類。濃厚系の活量・溶媒依存を全条件へ拡張しない。

\paragraph{問題の要旨} ハロゲン単体の状態・酸化力とHFの酸強度を比較する。

\Needspace{4\baselineskip}\paragraph{確認した条件・記述（独自要約）}\mbox{}\par

\noindent\textbf{判断対象：} ハロゲン単体の状態・酸化力とHFの酸強度を比較する。\par

\noindent\textbf{原典確認：} 22原典頁画像で五候補・否定・化学式を照合。\par

\paragraph{問われる核心} ハロゲン単体の状態・酸化力とHFの酸強度を比較する。

\paragraph{必要知識} 通常希薄水溶液ではHFは弱酸なので全強酸という候補5が不適切。

\paragraph{考え方・正答を選ぶ理由（査読待ち初稿）} 通常希薄水溶液ではHFは弱酸なので全強酸という候補5が不適切。F2単体の酸化力の強さ、HF水溶液の解離、毒性・腐食性は別の性質。酸化力系列と常温の状態を同じ五候補の順で照合する。

\Needspace{4\baselineskip}\paragraph{選択肢の検討（査読待ち初稿）}\mbox{}\par

\noindent\textbf{1：選ばない候補} 四単体は二原子分子で有毒。F2,Cl2,Br2,I2の分子単体と有害性。\par

\noindent\textbf{2：選ばない候補} F2/Cl2気体、Br2液体、I2固体。常温常圧の代表状態。\par

\noindent\textbf{3：選ばない候補} F2は水と激しく反応、I2は難溶。F2の酸化性とI2水への低溶解度。\par

\noindent\textbf{4：選ばない候補} 酸化力F2>Cl2>Br2>I2。標準的な水溶液の酸化還元比較。\par

\noindent\textbf{5：選ぶ候補} HF/HCl/HBr/HIは全強酸。通常希薄水溶液のHFは弱酸で、毒性と酸強度は別。\par

\paragraph{再利用する考え方} 原子番号と周期・族、物性と反応性、酸強度と危険性、生成と除去、定義の分母を別々に確かめる。

\paragraph{共通理論} \hyperref[theory:inorganic:halogen]{第\ref*{theory:inorganic:halogen}節}

\par 参考：\cite{inorganic_halogen_ref,niosh_hf}。

\paragraph{関連問題} \hyperref[q:r03-reg-applied-13]{R03-REG-APPLIED-13}, \hyperref[q:r07-reg-applied-17]{R07-REG-APPLIED-17}

\par 難易度 2（校準初稿）、優先度未評価。

\subsection{令和3年度 問13}
\label{q:r03-reg-applied-13}

\noindent\textbf{原典：} sources/R03/17\_applied\_science.pdf、PDF物理5頁。

\textbf{公式正答：} 1（single）。

\par\textbf{検討未確定：} 誘電率の全温度依存・Au/Ptの全混合物耐性は未照合。原液HFと水溶液濃度の条件を保持。

\paragraph{問題の要旨} HFの弱酸性と水素結合・腐食性を比較する。

\Needspace{4\baselineskip}\paragraph{確認した条件・記述（独自要約）}\mbox{}\par

\noindent\textbf{判断対象：} HFの弱酸性と水素結合・腐食性を比較する。\par

\noindent\textbf{原典確認：} 22原典頁画像で五候補・否定・化学式を照合。\par

\paragraph{問われる核心} HFの弱酸性と水素結合・腐食性を比較する。

\paragraph{必要知識} 通常HF水溶液を強酸とする1が不適切。

\paragraph{考え方・正答を選ぶ理由（査読待ち初稿）} 通常HF水溶液を強酸とする1が不適切。水中の解離が小さいことと腐食・有害性が低いことは同義でない。HF分子間水素結合はHX系列での高い沸点に関係する。NIOSHはHF液体/気体と40\%を越える水溶液の発煙を区別しており、原典の空気中発煙を全希薄水溶液の保証にはしない。

\Needspace{4\baselineskip}\paragraph{選択肢の検討（査読待ち初稿）}\mbox{}\par

\noindent\textbf{1：選ぶ候補} HF水溶液は毒性の強い強酸。毒性は強いが通常希薄水溶液は弱酸で酸強度の語が誤り。\par

\noindent\textbf{2：選ばない候補} 液体HFの高い誘電率。分極と液体の配向応答。値は温度等に依存。\par

\noindent\textbf{3：選ばない候補} 他HXより高沸点。分子間水素結合が異常な高沸点を説明。\par

\noindent\textbf{4：選ばない候補} 空気中で発煙。HF原液・濃厚水溶液で発煙。全希薄溶液とはしない。\par

\noindent\textbf{5：選ばない候補} SiO2/ガラスを腐食、Au/Ptを侵さない。ガラス除去性と貴金属の通常耐性。全混合条件保証ではない。\par

\paragraph{再利用する考え方} 原子番号と周期・族、物性と反応性、酸強度と危険性、生成と除去、定義の分母を別々に確かめる。

\paragraph{共通理論} \hyperref[theory:inorganic:halogen]{第\ref*{theory:inorganic:halogen}節}

\par 参考：\cite{niosh_hf,yale_hydrohalic}。

\paragraph{関連問題} \hyperref[q:h28-reg-applied-11]{H28-REG-APPLIED-11}

\par 難易度 2（校準初稿）、優先度未評価。

\subsection{令和7年度 問17}
\label{q:r07-reg-applied-17}

\noindent\textbf{原典：} sources/R07/17\_applied\_science.pdf、PDF物理6頁。

\textbf{公式正答：} 3（single）。

\par\textbf{検討未確定：} HFの酸強度は希薄水溶液の通常分類。濃厚系の活量・溶媒依存を全条件へ拡張しない。

\paragraph{問題の要旨} ハロゲン単体の状態・酸化力とHFの酸強度を比較する。

\Needspace{4\baselineskip}\paragraph{確認した条件・記述（独自要約）}\mbox{}\par

\noindent\textbf{判断対象：} ハロゲン単体の状態・酸化力とHFの酸強度を比較する。\par

\noindent\textbf{原典確認：} 22原典頁画像で五候補・否定・化学式を照合。\par

\paragraph{問われる核心} ハロゲン単体の状態・酸化力とHFの酸強度を比較する。

\paragraph{必要知識} 通常希薄水溶液ではHFは弱酸なので全強酸という候補3が不適切。

\paragraph{考え方・正答を選ぶ理由（査読待ち初稿）} 通常希薄水溶液ではHFは弱酸なので全強酸という候補3が不適切。F2単体の酸化力の強さ、HF水溶液の解離、毒性・腐食性は別の性質。酸化力系列と常温の状態を同じ五候補の順で照合する。

\Needspace{4\baselineskip}\paragraph{選択肢の検討（査読待ち初稿）}\mbox{}\par

\noindent\textbf{1：選ばない候補} F2/Cl2気体、Br2液体、I2固体。常温常圧の代表状態。\par

\noindent\textbf{2：選ばない候補} 酸化力F2>Cl2>Br2>I2。標準的な水溶液の酸化還元比較。\par

\noindent\textbf{3：選ぶ候補} HF/HCl/HBr/HIは全強酸。通常希薄水溶液のHFは弱酸で、毒性と酸強度は別。\par

\noindent\textbf{4：選ばない候補} F2は水と激しく反応、I2は難溶。F2の酸化性とI2水への低溶解度。\par

\noindent\textbf{5：選ばない候補} 四単体は二原子分子で有毒。F2,Cl2,Br2,I2の分子単体と有害性。\par

\paragraph{再利用する考え方} 原子番号と周期・族、物性と反応性、酸強度と危険性、生成と除去、定義の分母を別々に確かめる。

\paragraph{共通理論} \hyperref[theory:inorganic:halogen]{第\ref*{theory:inorganic:halogen}節}

\par 参考：\cite{inorganic_halogen_ref,niosh_hf}。

\paragraph{関連問題} \hyperref[q:h28-reg-applied-11]{H28-REG-APPLIED-11}

\par 難易度 2（校準初稿）、優先度未評価。

\section{水素化物の電荷・分極と多中心結合を比較する。}

\subsection{令和7年度 問16}
\label{q:r07-reg-applied-16}

\noindent\textbf{原典：} sources/R07/17\_applied\_science.pdf、PDF物理6頁。

\textbf{公式正答：} 5（single）。

\par\textbf{検討未確定：} B2H6最小安定の通常条件とBH3の特殊な単量体・配位付加体を区別。全遷移金属水素化物を不定比にしない。

\paragraph{問題の要旨} 水素化物の電荷・分極と多中心結合を比較する。

\Needspace{4\baselineskip}\paragraph{確認した条件・記述（独自要約）}\mbox{}\par

\noindent\textbf{判断対象：} 水素化物の電荷・分極と多中心結合を比較する。\par

\noindent\textbf{原典確認：} 22原典頁画像で五候補・否定・化学式を照合。\par

\paragraph{問われる核心} 水素化物の電荷・分極と多中心結合を比較する。

\paragraph{必要知識} 非金属のF/O等に結合したHは正に分極するので5が逆。

\paragraph{考え方・正答を選ぶ理由（査読待ち初稿）} 非金属のF/O等に結合したHは正に分極するので5が逆。金属塩型のH−、共有分子のHδ+、金属格子間隙の水素を区別する。B2H6の架橋B-H-Bには三中心二電子結合があり、全結合を普通の二中心結合だけで説明しない。最小安定という語はBH3の特殊条件での単量体存在を排除する普遍定義ではない。

\Needspace{4\baselineskip}\paragraph{選択肢の検討（査読待ち初稿）}\mbox{}\par

\noindent\textbf{1：選ばない候補} Hは外殻1/希ガスに1不足。H1s1はHe1s2へ1不足。Na等との共通点と独自性。\par

\noindent\textbf{2：選ばない候補} NaH/CaH2はイオン性。電気陽性金属とのH−を含む塩型。\par

\noindent\textbf{3：選ばない候補} 遷移金属間隙で不定比水素化物。格子の間隙占有率で組成が変わる系がある。\par

\noindent\textbf{4：選ばない候補} B2H6は最小安定ホウ素水素化物で3中心。通常持続する中性boraneの代表。BH3単量体の特殊条件存在と区別。\par

\noindent\textbf{5：選ぶ候補} ハロゲン/カルコゲンではHが負に分極。F/O等がHより電気陰性度大でHは通常正に分極、符号が逆。\par

\paragraph{再利用する考え方} 原子番号と周期・族、物性と反応性、酸強度と危険性、生成と除去、定義の分母を別々に確かめる。

\paragraph{共通理論} \hyperref[theory:inorganic:hydride]{第\ref*{theory:inorganic:hydride}節}

\par 参考：\cite{inorganic_hydrogen_ref,inorganic_metalloid_ref}。

\par 難易度 2（校準初稿）、優先度未評価。

\section{ゼオライトによる硬度イオン除去の機構を判断する。}

\subsection{平成23年度 問17}
\label{q:h23-reg-applied-17}

\noindent\textbf{原典：} sources/H23/17\_applied\_science.pdf、PDF物理6頁。

\textbf{公式正答：} 1（single）。

\paragraph{問題の要旨} ゼオライトによる硬度イオン除去の機構を判断する。

\Needspace{4\baselineskip}\paragraph{確認した条件・記述（独自要約）}\mbox{}\par

\noindent\textbf{判断対象：} ゼオライトによる硬度イオン除去の機構を判断する。\par

\noindent\textbf{原典確認：} 22原典頁画像で五候補・否定・化学式を照合。\par

\paragraph{問われる核心} ゼオライトによる硬度イオン除去の機構を判断する。

\paragraph{必要知識} ゼオライトの骨格負電荷を補償する交換可能陽イオンと、水中のCa2+/Mg2+が交換して硬度を下げるので1。

\paragraph{考え方・正答を選ぶ理由（査読待ち初稿）} ゼオライトの骨格負電荷を補償する交換可能陽イオンと、水中のCa2+/Mg2+が交換して硬度を下げるので1。2Na型+Ca2+⇌Ca型+2Na+という電荷収支で理解する。吸着機能もあることと、軟化に問われる主機構は別。

\Needspace{4\baselineskip}\paragraph{選択肢の検討（査読待ち初稿）}\mbox{}\par

\noindent\textbf{1：選ぶ候補} イオン交換。交換可能なNa等とCa/Mgを交換。\par

\noindent\textbf{2：選ばない候補} 吸着。表面への保持一般で、軟化の指定機構を十分に特定しない。\par

\noindent\textbf{3：選ばない候補} 吸収。内部への取込み一般で、対イオン交換を特定しない。\par

\noindent\textbf{4：選ばない候補} 還元。Ca/Mgの価数を減らす操作ではない。\par

\noindent\textbf{5：選ばない候補} 酸化。Ca/Mgを酸化する操作ではない。\par

\paragraph{再利用する考え方} 原子番号と周期・族、物性と反応性、酸強度と危険性、生成と除去、定義の分母を別々に確かめる。

\paragraph{共通理論} \hyperref[theory:inorganic:exchange]{第\ref*{theory:inorganic:exchange}節}

\par 参考：\cite{inorganic_zeolite}。

\par 難易度 2（校準初稿）、優先度未評価。

\section{希ガスの通常状態・外殻例外と放電を比較する。}

\subsection{平成23年度 問11}
\label{q:h23-reg-applied-11}

\noindent\textbf{原典：} sources/H23/17\_applied\_science.pdf、PDF物理4頁。

\textbf{公式正答：} 4（single）。

\paragraph{問題の要旨} 希ガスの通常状態・外殻例外と放電を比較する。

\Needspace{4\baselineskip}\paragraph{確認した条件・記述（独自要約）}\mbox{}\par

\noindent\textbf{判断対象：} 希ガスの通常状態・外殻例外と放電を比較する。\par

\noindent\textbf{原典確認：} 22原典頁画像で五候補・否定・化学式を照合。\par

\paragraph{問われる核心} 希ガスの通常状態・外殻例外と放電を比較する。

\paragraph{必要知識} 希ガス原子の励起・脱励起が特有の放電色を与えるので4。

\paragraph{考え方・正答を選ぶ理由（査読待ち初稿）} 希ガス原子の励起・脱励起が特有の放電色を与えるので4。通常の空気中は単体、常温常圧で気体でありHe外殻は2。Xe化合物の存在から希ガスは一切反応しないとする逆の全称にも広げない。沸点の低さと化学的反応性は別の性質。

\Needspace{4\baselineskip}\paragraph{選択肢の検討（査読待ち初稿）}\mbox{}\par

\noindent\textbf{1：選ばない候補} 空気中は化合物。通常の空気では単体原子で存在。\par

\noindent\textbf{2：選ばない候補} 常温常圧で気体と液体。出題範囲のHeからRnは通常気体。\par

\noindent\textbf{3：選ばない候補} 軽いほど低沸点・高反応。沸点傾向はよいが軽いほど反応性に富むは適切でない。\par

\noindent\textbf{4：選ぶ候補} 低圧放電で特有色。励起原子の準位遷移の発光。\par

\noindent\textbf{5：選ばない候補} 全て外殻8電子。Heの1s2は外殻2電子という例外。\par

\paragraph{再利用する考え方} 原子番号と周期・族、物性と反応性、酸強度と危険性、生成と除去、定義の分母を別々に確かめる。

\paragraph{共通理論} \hyperref[theory:inorganic:noble]{第\ref*{theory:inorganic:noble}節}

\par 参考：\cite{inorganic_noble_ref}。

\paragraph{関連問題} \hyperref[q:h27-reg-applied-10]{H27-REG-APPLIED-10}

\par 難易度 2（校準初稿）、優先度未評価。

\subsection{平成27年度 問10}
\label{q:h27-reg-applied-10}

\noindent\textbf{原典：} sources/H27/17\_applied\_science.pdf、PDF物理4頁。

\textbf{公式正答：} 4（single）。

\paragraph{問題の要旨} 希ガスの通常状態・外殻例外と放電を比較する。

\Needspace{4\baselineskip}\paragraph{確認した条件・記述（独自要約）}\mbox{}\par

\noindent\textbf{判断対象：} 希ガスの通常状態・外殻例外と放電を比較する。\par

\noindent\textbf{原典確認：} 22原典頁画像で五候補・否定・化学式を照合。\par

\paragraph{問われる核心} 希ガスの通常状態・外殻例外と放電を比較する。

\paragraph{必要知識} 希ガス原子の励起・脱励起が特有の放電色を与えるので4。

\paragraph{考え方・正答を選ぶ理由（査読待ち初稿）} 希ガス原子の励起・脱励起が特有の放電色を与えるので4。通常の空気中は単体、常温常圧で気体でありHe外殻は2。Xe化合物の存在から希ガスは一切反応しないとする逆の全称にも広げない。沸点の低さと化学的反応性は別の性質。

\Needspace{4\baselineskip}\paragraph{選択肢の検討（査読待ち初稿）}\mbox{}\par

\noindent\textbf{1：選ばない候補} 空気中は化合物。通常の空気では単体原子で存在。\par

\noindent\textbf{2：選ばない候補} 常温常圧で気体と液体。出題範囲のHeからRnは通常気体。\par

\noindent\textbf{3：選ばない候補} 軽いほど低沸点・高反応。沸点傾向はよいが軽いほど反応性に富むは適切でない。\par

\noindent\textbf{4：選ぶ候補} 低圧放電で特有色。励起原子の準位遷移の発光。\par

\noindent\textbf{5：選ばない候補} 全て外殻8電子。Heの1s2は外殻2電子という例外。\par

\paragraph{再利用する考え方} 原子番号と周期・族、物性と反応性、酸強度と危険性、生成と除去、定義の分母を別々に確かめる。

\paragraph{共通理論} \hyperref[theory:inorganic:noble]{第\ref*{theory:inorganic:noble}節}

\par 参考：\cite{inorganic_noble_ref}。

\paragraph{関連問題} \hyperref[q:h23-reg-applied-11]{H23-REG-APPLIED-11}, \hyperref[q:r01-ret-applied-16]{R01-RET-APPLIED-16}

\par 難易度 2（校準初稿）、優先度未評価。

\subsection{令和元年度再試験 問16}
\label{q:r01-ret-applied-16}

\noindent\textbf{原典：} sources/R01\_retest/17\_applied\_science.pdf、PDF物理6頁。

\textbf{公式正答：} 2（single）。

\paragraph{問題の要旨} 希ガスの単原子性とHeの電子例外を確認する。

\Needspace{4\baselineskip}\paragraph{確認した条件・記述（独自要約）}\mbox{}\par

\noindent\textbf{判断対象：} 希ガスの単原子性とHeの電子例外を確認する。\par

\noindent\textbf{原典確認：} 22原典頁画像で五候補・否定・化学式を照合。\par

\paragraph{問われる核心} 希ガスの単原子性とHeの電子例外を確認する。

\paragraph{必要知識} 原典の1原子分子は単原子気体という意味で2を選ぶ。

\paragraph{考え方・正答を選ぶ理由（査読待ち初稿）} 原典の1原子分子は単原子気体という意味で2を選ぶ。共有結合した多原子分子ではない。Heの1s2とXeの化合物の例外を保持し、通常状態と極低温や特殊な分子種を混同しない。

\Needspace{4\baselineskip}\paragraph{選択肢の検討（査読待ち初稿）}\mbox{}\par

\noindent\textbf{1：選ばない候補} 空気中は化合物。通常空気は単体。\par

\noindent\textbf{2：選ぶ候補} 1原子分子からなる。通常の単体気体は単原子で2が適切。\par

\noindent\textbf{3：選ばない候補} 常温常圧で気体と液体。出題範囲の通常単体は気体。\par

\noindent\textbf{4：選ばない候補} 軽いほど低沸点・高反応。軽いほど高反応性とは言えない。\par

\noindent\textbf{5：選ばない候補} 全て外殻8電子。Heは2で例外。\par

\paragraph{再利用する考え方} 原子番号と周期・族、物性と反応性、酸強度と危険性、生成と除去、定義の分母を別々に確かめる。

\paragraph{共通理論} \hyperref[theory:inorganic:noble]{第\ref*{theory:inorganic:noble}節}

\par 参考：\cite{inorganic_noble_ref}。

\paragraph{関連問題} \hyperref[q:h27-reg-applied-10]{H27-REG-APPLIED-10}

\par 難易度 2（校準初稿）、優先度未評価。

\section{アルカリ金属の半径・イオン化と融点を比較する。}

\subsection{平成24年度 問12}
\label{q:h24-reg-applied-12}

\noindent\textbf{原典：} sources/H24/17\_applied\_science.pdf、PDF物理4頁。

\textbf{公式正答：} 4（single）。

\paragraph{問題の要旨} アルカリ金属の半径・イオン化と融点を比較する。

\Needspace{4\baselineskip}\paragraph{確認した条件・記述（独自要約）}\mbox{}\par

\noindent\textbf{判断対象：} アルカリ金属の半径・イオン化と融点を比較する。\par

\noindent\textbf{原典確認：} 22原典頁画像で五候補・否定・化学式を照合。\par

\paragraph{問われる核心} アルカリ金属の半径・イオン化と融点を比較する。

\paragraph{必要知識} Li,Na,K,Rb,Csは外殻ns1の第1族で、殻が増え陽イオン半径が大きく第一IEは低くなる。

\paragraph{考え方・正答を選ぶ理由（査読待ち初稿）} Li,Na,K,Rb,Csは外殻ns1の第1族で、殻が増え陽イオン半径が大きく第一IEは低くなる。通常圧の融点はLiからCsへ下がるので4が不適切。融点は金属結晶の凝集と相の自由エネルギーに関係し、原子のIEの同義ではない。

\Needspace{4\baselineskip}\paragraph{選択肢の検討（査読待ち初稿）}\mbox{}\par

\noindent\textbf{1：選ばない候補} 外殻s1が失われ希ガス配置。電子を1個失う通常の1価陽イオン。\par

\noindent\textbf{2：選ばない候補} Li→Csで陽イオン半径増。同じ電荷で殻が増える系列。\par

\noindent\textbf{3：選ばない候補} Li→CsでIE低下。外殻が遠く遮蔽されて第一IEが下がる。\par

\noindent\textbf{4：選ぶ候補} Li→Csで融点上昇。通常の融点は系列で低下するので逆。\par

\noindent\textbf{5：選ばない候補} 固有の炎色。原子励起の発光スペクトルの違い。\par

\paragraph{再利用する考え方} 原子番号と周期・族、物性と反応性、酸強度と危険性、生成と除去、定義の分母を別々に確かめる。

\paragraph{共通理論} \hyperref[theory:inorganic:periodic]{第\ref*{theory:inorganic:periodic}節}

\par 参考：\cite{inorganic_alkali,atomic_nist2019}。

\par 難易度 2（校準初稿）、優先度未評価。

\subsection{平成29年度 問8}
\label{q:h29-reg-applied-08}

\noindent\textbf{原典：} sources/H29/17\_applied\_science.pdf、PDF物理3頁。

\textbf{公式正答：} 5（single）。

\paragraph{問題の要旨} 周期表の副殻容量と希土類の電子配置を比較する。

\Needspace{4\baselineskip}\paragraph{確認した条件・記述（独自要約）}\mbox{}\par

\noindent\textbf{判断対象：} 周期表の副殻容量と希土類の電子配置を比較する。\par

\noindent\textbf{原典確認：} 22原典頁画像で五候補・否定・化学式を照合。\par

\paragraph{問われる核心} 周期表の副殻容量と希土類の電子配置を比較する。

\paragraph{必要知識} 希土類をd副殻に1個ずつ増える系列とする5は誤り。

\paragraph{考え方・正答を選ぶ理由（査読待ち初稿）} 希土類をd副殻に1個ずつ増える系列とする5は誤り。ランタノイドのf電子数にも例外があり、希土類の定義はSc/Yを含む。閉副殻Beと閉殻希ガスの化学性も同一視しない。原子番号・副殻・族の三つの分類を分ける。

\Needspace{4\baselineskip}\paragraph{選択肢の検討（査読待ち初稿）}\mbox{}\par

\noindent\textbf{1：選ばない候補} 第3周期はn3まで占有。中性基底の第3周期の配置。\par

\noindent\textbf{2：選ばない候補} s/p/dは2/6/10電子。副殻軌道数1/3/5と各2電子。\par

\noindent\textbf{3：選ばない候補} Beの2s満でも不活性ではない。閉じた2s副殻だけで希ガス的全不活性にならない。\par

\noindent\textbf{4：選ばない候補} K19番目は4s。通常K基底[Ar]4s1。\par

\noindent\textbf{5：選ぶ候補} 希土類はd電子が1ずつ増える。主にf配置が変わる系列で、dの単調増加ではない。\par

\paragraph{再利用する考え方} 原子番号と周期・族、物性と反応性、酸強度と危険性、生成と除去、定義の分母を別々に確かめる。

\paragraph{共通理論} \hyperref[theory:inorganic:periodic]{第\ref*{theory:inorganic:periodic}節}

\par 参考：\cite{atomic_nist2019,inorganic_periodic_ref}。

\paragraph{関連問題} \hyperref[q:h30-reg-applied-10]{H30-REG-APPLIED-10}

\par 難易度 2（校準初稿）、優先度未評価。

\subsection{平成30年度 問10}
\label{q:h30-reg-applied-10}

\noindent\textbf{原典：} sources/H30/17\_applied\_science.pdf、PDF物理3頁。

\textbf{公式正答：} 3（single）。

\par\textbf{検討未確定：} 主族/遷移元素の定義には12族の取扱い差。dblockの分類と区別する。

\paragraph{問題の要旨} 電子配置の例外と周期表ブロックの定義を比較する。

\Needspace{4\baselineskip}\paragraph{確認した条件・記述（独自要約）}\mbox{}\par

\noindent\textbf{判断対象：} 電子配置の例外と周期表ブロックの定義を比較する。\par

\noindent\textbf{原典確認：} 22原典頁画像で五候補・否定・化学式を照合。\par

\paragraph{問われる核心} 電子配置の例外と周期表ブロックの定義を比較する。

\paragraph{必要知識} CuとZnの基底配置の4s電子数だけを見てsblockと結論する3は誤り。

\paragraph{考え方・正答を選ぶ理由（査読待ち初稿）} CuとZnの基底配置の4s電子数だけを見てsblockと結論する3は誤り。ブロックは周期表のd系列の位置として分類し、個別基底配置の例外だけで移動させない。遷移元素の語では12族を除外する定義があり、dblockとの同義を強制しない。

\Needspace{4\baselineskip}\paragraph{選択肢の検討（査読待ち初稿）}\mbox{}\par

\noindent\textbf{1：選ばない候補} s/pは主族。通常のs/p主族分類。\par

\noindent\textbf{2：選ばない候補} dは遷移、12族を除く場合あり。遷移元素の定義とdblockの範囲を区別。\par

\noindent\textbf{3：選ぶ候補} Cu4s1d10/Zn4s2d10だからsblock。電子配置は正しくてもCu/Znの周期表位置はdblock。\par

\noindent\textbf{4：選ばない候補} lanthanoid/actinoid内部遷移。内部系列の代表分類。\par

\noindent\textbf{5：選ばない候補} fblockは1原子番号で必ずf1増でない。基底配置の例外を認める。\par

\paragraph{再利用する考え方} 原子番号と周期・族、物性と反応性、酸強度と危険性、生成と除去、定義の分母を別々に確かめる。

\paragraph{共通理論} \hyperref[theory:inorganic:periodic]{第\ref*{theory:inorganic:periodic}節}

\par 参考：\cite{atomic_nist2019}。

\paragraph{関連問題} \hyperref[q:h29-reg-applied-08]{H29-REG-APPLIED-08}

\par 難易度 2（校準初稿）、優先度未評価。

\subsection{令和5年度 問12}
\label{q:r05-reg-applied-12}

\noindent\textbf{原典：} sources/R05/17\_applied\_science.pdf、PDF物理5頁。

\textbf{公式正答：} 5（single）。

\paragraph{問題の要旨} 元素の原子番号を周期表で比較する。

\Needspace{4\baselineskip}\paragraph{確認した条件・記述（独自要約）}\mbox{}\par

\noindent\textbf{判断対象：} 元素の原子番号を周期表で比較する。\par

\noindent\textbf{原典確認：} 22原典頁画像で五候補・否定・化学式を照合。\par

\paragraph{問われる核心} 元素の原子番号を周期表で比較する。

\paragraph{必要知識} 原子番号は陽子数で質量順位ではない。

\paragraph{考え方・正答を選ぶ理由（査読待ち初稿）} 原子番号は陽子数で質量順位ではない。Nb41/Ga31、Pm61/Dy66、B5/P15、I53/Kr36の比較は正しいがRe75はTa73より大きく5が不適切。原子質量や天然存在量の大小から推測しない。

\Needspace{4\baselineskip}\paragraph{選択肢の検討（査読待ち初稿）}\mbox{}\par

\noindent\textbf{1：選ばない候補} Nb>Ga。41>31。\par

\noindent\textbf{2：選ばない候補} Pm<Dy。61<66。\par

\noindent\textbf{3：選ばない候補} B<P。5<15。\par

\noindent\textbf{4：選ばない候補} I>Kr。53>36。\par

\noindent\textbf{5：選ぶ候補} Re<Ta。Re75>Ta73で小さいは逆。\par

\paragraph{再利用する考え方} 原子番号と周期・族、物性と反応性、酸強度と危険性、生成と除去、定義の分母を別々に確かめる。

\paragraph{共通理論} \hyperref[theory:inorganic:periodic]{第\ref*{theory:inorganic:periodic}節}

\par 参考：\cite{atomic_nist2019}。

\paragraph{関連問題} \hyperref[q:r06-reg-applied-13]{R06-REG-APPLIED-13}

\par 難易度 2（校準初稿）、優先度未評価。

\subsection{令和6年度 問13}
\label{q:r06-reg-applied-13}

\noindent\textbf{原典：} sources/R06/17\_applied\_science.pdf、PDF物理4頁。

\textbf{公式正答：} 5（single）。

\paragraph{問題の要旨} 周期と族を分けて元素の位置を比較する。

\Needspace{4\baselineskip}\paragraph{確認した条件・記述（独自要約）}\mbox{}\par

\noindent\textbf{判断対象：} 周期と族を分けて元素の位置を比較する。\par

\noindent\textbf{原典確認：} 22原典頁画像で五候補・否定・化学式を照合。\par

\paragraph{問われる核心} 周期と族を分けて元素の位置を比較する。

\paragraph{必要知識} MoとIrの周期を同じ6とする5が不適切。

\paragraph{考え方・正答を選ぶ理由（査読待ち初稿）} MoとIrの周期を同じ6とする5が不適切。周期は横の列、族は縦の列。Mg/Si第3、In/B13族、Nb/Tc第5、Au/Cu11族の五対応を元素の位置から確認する。

\Needspace{4\baselineskip}\paragraph{選択肢の検討（査読待ち初稿）}\mbox{}\par

\noindent\textbf{1：選ばない候補} Mg/Siは第3周期。Mg12/Si14は第3。\par

\noindent\textbf{2：選ばない候補} In/Bは同族。In49/B5は13族。\par

\noindent\textbf{3：選ばない候補} Nb/Tcは第5周期。Nb41/Tc43は第5。\par

\noindent\textbf{4：選ばない候補} Au/Cuは同族。Au79/Cu29は11族。\par

\noindent\textbf{5：選ぶ候補} Mo/Irは第6周期。Mo42は第5、Ir77は第6。\par

\paragraph{再利用する考え方} 原子番号と周期・族、物性と反応性、酸強度と危険性、生成と除去、定義の分母を別々に確かめる。

\paragraph{共通理論} \hyperref[theory:inorganic:periodic]{第\ref*{theory:inorganic:periodic}節}

\par 参考：\cite{atomic_nist2019}。

\paragraph{関連問題} \hyperref[q:r05-reg-applied-12]{R05-REG-APPLIED-12}

\par 難易度 2（校準初稿）、優先度未評価。

\subsection{令和7年度 問15}
\label{q:r07-reg-applied-15}

\noindent\textbf{原典：} sources/R07/17\_applied\_science.pdf、PDF物理5頁。

\textbf{公式正答：} 2（single）。

\paragraph{問題の要旨} ランタノイドとアクチノイドを区別する。

\Needspace{4\baselineskip}\paragraph{確認した条件・記述（独自要約）}\mbox{}\par

\noindent\textbf{判断対象：} ランタノイドとアクチノイドを区別する。\par

\noindent\textbf{原典確認：} 22原典頁画像で五候補・否定・化学式を照合。\par

\paragraph{問われる核心} ランタノイドとアクチノイドを区別する。

\paragraph{必要知識} Th90はアクチノイドなので2。

\paragraph{考え方・正答を選ぶ理由（査読待ち初稿）} Th90はアクチノイドなので2。Sm62,Nd60,Pm61,Ho67はランタノイド。放射性か安定かという軸は系列分類と別で、放射性Pmを系列外としない。希土類の広い範囲とランタノイド15元素も区別する。

\Needspace{4\baselineskip}\paragraph{選択肢の検討（査読待ち初稿）}\mbox{}\par

\noindent\textbf{1：選ばない候補} Sm。Sm62はlanthanoid。\par

\noindent\textbf{2：選ぶ候補} Th。Th90はactinoidで提示系列に属さない。\par

\noindent\textbf{3：選ばない候補} Nd。Nd60はlanthanoid。\par

\noindent\textbf{4：選ばない候補} Pm。Pm61は放射性だがlanthanoid。\par

\noindent\textbf{5：選ばない候補} Ho。Ho67はlanthanoid。\par

\paragraph{再利用する考え方} 原子番号と周期・族、物性と反応性、酸強度と危険性、生成と除去、定義の分母を別々に確かめる。

\paragraph{共通理論} \hyperref[theory:inorganic:periodic]{第\ref*{theory:inorganic:periodic}節}

\par 参考：\cite{atomic_nist2019}。

\par 難易度 2（校準初稿）、優先度未評価。

\section{金属イオンの沈殿・錯体と反応条件を確認する。}

\subsection{平成25年度 問13}
\label{q:h25-reg-applied-13}

\noindent\textbf{原典：} sources/H25/17\_applied\_science.pdf、PDF物理6頁。

\textbf{公式正答：} 4（single）。

\par\textbf{検討未確定：} 酸性ZnSの析出はpH・濃度に依存し3の一律排除をしない。濃硝酸Feの温度・濃度・表面状態も未指定。

\paragraph{問題の要旨} 金属イオンの沈殿・錯体と反応条件を確認する。

\Needspace{4\baselineskip}\paragraph{確認した条件・記述（独自要約）}\mbox{}\par

\noindent\textbf{判断対象：} 金属イオンの沈殿・錯体と反応条件を確認する。\par

\noindent\textbf{原典確認：} 22原典頁画像で五候補・否定・化学式を照合。\par

\paragraph{問われる核心} 金属イオンの沈殿・錯体と反応条件を確認する。

\paragraph{必要知識} NH3水は塩基としてOH−を供給してZn(OH)2を析出させ、過剰NH3は[Zn(NH3)4]2+として溶解を促すので4。

\paragraph{考え方・正答を選ぶ理由（査読待ち初稿）} NH3水は塩基としてOH−を供給してZn(OH)2を析出させ、過剰NH3は[Zn(NH3)4]2+として溶解を促すので4。色の誤り、光分解と溶解、酸の酸化性と不動態を分ける。3の酸性という広い表現はZnSが全酸性で析出不可という証明にはならず、pH・金属濃度・S2−の活量に留保を残す。

\Needspace{4\baselineskip}\paragraph{選択肢の検討（査読待ち初稿）}\mbox{}\par

\noindent\textbf{1：選ばない候補} BaSO4黒色沈殿。BaSO4の通常沈殿は白色。\par

\noindent\textbf{2：選ばない候補} AgCl白沈殿が日光で溶解。光分解で銀等を生じ暗色化し、溶解とは異なる。\par

\noindent\textbf{3：選ばない候補} 酸性ZnにH2Sで白沈殿。ZnS白色はあり得るが、酸性の程度・濃度で析出が変わる。\par

\noindent\textbf{4：選ぶ候補} Znに少量NH3で白沈殿/過剰で溶解。Zn(OH)2沈殿と過剰NH3のammine錯体で溶解。\par

\noindent\textbf{5：選ばない候補} Fe濃硝酸でH2と溶解。通常濃硝酸はFeを不動態化。酸化性酸からH2発生とも説明できない。\par

\paragraph{再利用する考え方} 原子番号と周期・族、物性と反応性、酸強度と危険性、生成と除去、定義の分母を別々に確かめる。

\paragraph{共通理論} \hyperref[theory:inorganic:complex]{第\ref*{theory:inorganic:complex}節}

\par 参考：\cite{qualitative_stanford,inorganic_zinc_ref}。

\par 難易度 2（校準初稿）、優先度未評価。

\section{Siの通常性状と炭化物・ホウ化物形成を比較する。}

\subsection{令和2年度 問14}
\label{q:r02-reg-applied-14}

\noindent\textbf{原典：} sources/R02/17\_applied\_science.pdf、PDF物理5頁。

\textbf{公式正答：} 4（single）。

\par\textbf{検討未確定：} 水/無機酸不溶の全称には混酸等の例外。純度・酸種・温度・表面条件未指定。公式誤りと他候補の一般化疑義を分ける。

\paragraph{問題の要旨} Siの通常性状と炭化物・ホウ化物形成を比較する。

\Needspace{4\baselineskip}\paragraph{確認した条件・記述（独自要約）}\mbox{}\par

\noindent\textbf{判断対象：} Siの通常性状と炭化物・ホウ化物形成を比較する。\par

\noindent\textbf{原典確認：} 22原典頁画像で五候補・否定・化学式を照合。\par

\paragraph{問われる核心} Siの通常性状と炭化物・ホウ化物形成を比較する。

\paragraph{必要知識} SiCやSiのホウ化物が存在するため、C/Bと結合しないという候補4が不適切。

\paragraph{考え方・正答を選ぶ理由（査読待ち初稿）} SiCやSiのホウ化物が存在するため、C/Bと結合しないという候補4が不適切。半導体用途、地殻での酸化物/ケイ酸塩、diamond型構造を区別する。一方、無機酸に全く溶けないという他候補は酸種・混合・酸化剤で変わり得るので留保を残す。

\Needspace{4\baselineskip}\paragraph{選択肢の検討（査読待ち初稿）}\mbox{}\par

\noindent\textbf{1：選ばない候補} 高純度Siは半導体材料。通常の電子材料用途。\par

\noindent\textbf{2：選ばない候補} SiO2/ケイ酸塩で地殻に多い。元素Siの単体ではなく化合物の分布。\par

\noindent\textbf{3：選ばない候補} 高純度Siは水/無機酸に不溶。通常の非酸化性酸の代表耐性。酸化剤とHFを組合せたエッチングなどに例外。\par

\noindent\textbf{4：選ぶ候補} 高温でもC/Bと直接結合しない。SiCとSiB6等の化合物が存在し無結合の全称は誤り。\par

\noindent\textbf{5：選ばない候補} Siにはdiamond構造。通常結晶Siは四面体結合のdiamond型。\par

\paragraph{再利用する考え方} 原子番号と周期・族、物性と反応性、酸強度と危険性、生成と除去、定義の分母を別々に確かめる。

\paragraph{共通理論} \hyperref[theory:inorganic:silicon]{第\ref*{theory:inorganic:silicon}節}

\par 参考：\cite{inorganic_metalloid_ref,inorganic_sic,inorganic_sib6,inorganic_hf_etch}。

\paragraph{関連問題} \hyperref[q:r06-reg-applied-09]{R06-REG-APPLIED-09}

\par 難易度 2（校準初稿）、優先度未評価。

\subsection{令和6年度 問9}
\label{q:r06-reg-applied-09}

\noindent\textbf{原典：} sources/R06/17\_applied\_science.pdf、PDF物理3頁。

\textbf{公式正答：} 5（single）。

\par\textbf{検討未確定：} 水/無機酸不溶の全称には混酸等の例外。純度・酸種・温度・表面条件未指定。公式誤りと他候補の一般化疑義を分ける。

\paragraph{問題の要旨} Siの通常性状と炭化物・ホウ化物形成を比較する。

\Needspace{4\baselineskip}\paragraph{確認した条件・記述（独自要約）}\mbox{}\par

\noindent\textbf{判断対象：} Siの通常性状と炭化物・ホウ化物形成を比較する。\par

\noindent\textbf{原典確認：} 22原典頁画像で五候補・否定・化学式を照合。\par

\paragraph{問われる核心} Siの通常性状と炭化物・ホウ化物形成を比較する。

\paragraph{必要知識} SiCやSiのホウ化物が存在するため、C/Bと結合しないという候補5が不適切。

\paragraph{考え方・正答を選ぶ理由（査読待ち初稿）} SiCやSiのホウ化物が存在するため、C/Bと結合しないという候補5が不適切。半導体用途、地殻での酸化物/ケイ酸塩、diamond型構造を区別する。一方、無機酸に全く溶けないという他候補は酸種・混合・酸化剤で変わり得るので留保を残す。

\Needspace{4\baselineskip}\paragraph{選択肢の検討（査読待ち初稿）}\mbox{}\par

\noindent\textbf{1：選ばない候補} SiO2/ケイ酸塩で地殻に多い。元素Siの単体ではなく化合物の分布。\par

\noindent\textbf{2：選ばない候補} 高純度Siは半導体材料。通常の電子材料用途。\par

\noindent\textbf{3：選ばない候補} 高純度Siは水/無機酸に不溶。通常の非酸化性酸の代表耐性。酸化剤とHFを組合せたエッチングなどに例外。\par

\noindent\textbf{4：選ばない候補} Siにはdiamond構造。通常結晶Siは四面体結合のdiamond型。\par

\noindent\textbf{5：選ぶ候補} 高温でもC/Bと直接結合しない。SiCとSiB6等の化合物が存在し無結合の全称は誤り。\par

\paragraph{再利用する考え方} 原子番号と周期・族、物性と反応性、酸強度と危険性、生成と除去、定義の分母を別々に確かめる。

\paragraph{共通理論} \hyperref[theory:inorganic:silicon]{第\ref*{theory:inorganic:silicon}節}

\par 参考：\cite{inorganic_metalloid_ref,inorganic_sic,inorganic_sib6,inorganic_hf_etch}。

\paragraph{関連問題} \hyperref[q:r02-reg-applied-14]{R02-REG-APPLIED-14}

\par 難易度 2（校準初稿）、優先度未評価。

\section{分子・結晶・錯体の配位を比較する。}

\subsection{平成29年度 問10}
\label{q:h29-reg-applied-10}

\noindent\textbf{原典：} sources/H29/17\_applied\_science.pdf、PDF物理3頁。

\textbf{公式正答：} 4（single）。

\par\textbf{検討未確定：} ホウ素全同素体の一般化とヘムの厳密理想配位には留保。原典の五配位/六配位と分子全構造を区別。

\paragraph{問題の要旨} 分子・結晶・錯体の配位を比較する。

\Needspace{4\baselineskip}\paragraph{確認した条件・記述（独自要約）}\mbox{}\par

\noindent\textbf{判断対象：} 分子・結晶・錯体の配位を比較する。\par

\noindent\textbf{原典確認：} 22原典頁画像で五候補・否定・化学式を照合。\par

\paragraph{問われる核心} 分子・結晶・錯体の配位を比較する。

\paragraph{必要知識} NaClとCsClは異なる構造で4が不適切。

\paragraph{考え方・正答を選ぶ理由（査読待ち初稿）} NaClとCsClは異なる構造で4が不適切。NaClはFCC基底の6配位、CsClは単純立方基底の8配位。P4とBXのクラスター、四面体分子、ヘムの配位数は別の階層。ヘモグロビンの近似配位を完全な理想八面体全角度の保証にしない。

\Needspace{4\baselineskip}\paragraph{選択肢の検討（査読待ち初稿）}\mbox{}\par

\noindent\textbf{1：選ばない候補} 白リンP4は正四面体。P4の4原子の分子構造。\par

\noindent\textbf{2：選ばない候補} Bは12頂点二十面体単位。代表的な結晶BのB12単位。全同素体の例外は留保。\par

\noindent\textbf{3：選ばない候補} CCl4/SiCl4は四面体。中心原子に4結合の分子形。\par

\noindent\textbf{4：選ぶ候補} NaCl/CsClは同一結晶型。NaCl6/CsCl8の配位とBravaisが違う。\par

\noindent\textbf{5：選ばない候補} deoxyheme5配位、oxyheme6配位。酸素結合で第6配位子が加わる代表近似。平面からの変位等もある。\par

\paragraph{再利用する考え方} 原子番号と周期・族、物性と反応性、酸強度と危険性、生成と除去、定義の分母を別々に確かめる。

\paragraph{共通理論} \hyperref[theory:inorganic:structure]{第\ref*{theory:inorganic:structure}節}

\par 参考：\cite{openstax_crystals,inorganic_metalloid_ref,inorganic_heme}。

\par 難易度 2（校準初稿）、優先度未評価。
